\documentclass[11pt]{article}
\usepackage[margin=1in]{geometry}
\usepackage[T1]{fontenc}
\usepackage{lmodern}
\usepackage{hyperref}
\usepackage{booktabs}

\title{A Portable Finite-State Architecture for Deterministic Checkers Play}
\author{Christopher A. Tucker\\
\texttt{cartheur@pm.me}}
\date{Draft --- November 2026}

\begin{document}
\maketitle

\begin{abstract}
This paper introduces a work-in-progress architecture for a deterministic
Checkers player designed to run across constrained and desktop environments.
The controller is expressed as a finite-state machine (FSA) with explicit
request, validation, rule generation, bounded projection, evaluation,
selection, response, and error states. The design separates the control
machine from policy: the FSA governs safe, inspectable progress, while a
replaceable policy ranks legal candidate moves. A portable C++11 core, a
word-oriented mailbox protocol, board-geometry fixtures, and initial
men-only rules fixtures provide the current implementation evidence. This
draft reports the architecture and its validation approach; it does not yet
claim game-playing strength, complete Checkers rules, or novelty relative to
prior game programs.
\end{abstract}

\section{Introduction}

Classic game systems demonstrate that useful computer players need not rely on
large learned models. This project investigates a different question: can a
small, deterministic controller provide a transparent and portable foundation
for Checkers play across a host emulator and constrained target hardware?

The proposed answer is an FSA with explicit state, bounded work, and observable
mailbox results. It is intended to make each stage inspectable and
fixture-testable rather than embedding control flow implicitly in a monolithic
search routine.

\section{Architecture}

The controller begins in $q_0=\mathrm{WAIT\_REQUEST}$. It snapshots and
validates an incoming board, reconciles prior results, generates legal moves,
and enters $q_4=\mathrm{ANTICIPATE}$. In that state it performs bounded
projection of candidate lines before evaluation and selection. A retained
move enters $q_7=\mathrm{DECISION}$, whose transition writes the successful
response; invalid or terminal work enters $q_E=\mathrm{ERROR}$, whose
transition writes the error response. The current state diagram and formal
input/output alphabets are maintained with the source repository.

\begin{center}
\fbox{\parbox{0.9\linewidth}{
\textbf{Current evidence artifact:} \texttt{research/checkers-evidence/FSA-LOGIC.md}.
The rendered diagram is \texttt{FSA-LOGIC.png}. At this draft stage the
portable core implements the request/response skeleton; move generation
provides evidence for the next generation state.
}}
\end{center}

\subsection{Separation of control and policy}

The FSA is not the game strategy. It defines state transitions, memory
ownership, failure paths, and commit boundaries. In particular,
\texttt{ANTICIPATE} is a controller state, while projection is the bounded
operation it performs; \texttt{DECISION} and \texttt{ERROR} name conditions,
while their outgoing transitions perform mailbox writes. Policy is a separate,
deterministic ranking function over legal projected moves. This separation
permits policy experiments without weakening protocol or rules validation.

\section{Anticipation and decision}

The central distinction in this architecture is between anticipating a
consequence and deciding what to expose to the game surface. The distinction
is useful even for a small deterministic player because it places board
simulation, policy, memory, and external effects on different sides of an
explicit boundary.

\subsection{\texttt{ANTICIPATE}: bounded projection, not commitment}

The \texttt{ANTICIPATE} state receives a legal candidate from generation and
constructs a private projected position. At minimum, projection applies the
candidate move and records captures, promotion, turn ownership, and any
required jump continuation. A later implementation may extend the same state
to project a bounded opponent reply and a reply-to-reply line. Thus,
anticipation is not a prediction that the opponent must make a particular
move; it is a finite, reproducible construction of positions that may result
from a candidate line.

Three constraints make this state suitable for a constrained system:
\begin{enumerate}
  \item \textbf{Legality precedes projection.} The state consumes only a
  generated legal move; it does not invent board transitions or bypass the
  rules engine.
  \item \textbf{Projection is private and bounded.} The accepted board remains
  immutable while a separate scratch position and transition trace are built.
  Each FSA invocation performs a bounded unit of work, such as applying one
  candidate or one continuation, before yielding control.
  \item \textbf{Projection remains inspectable.} The trace retains more than a
  numerical value: it identifies the candidate, projected board, captures,
  continuations, and later policy value. It can therefore explain why a
  selected move was possible and support reconciliation with the next observed
  board.
\end{enumerate}

The state may loop while further candidate lines need projection. Its exit
condition is not a likely move being found, but that the explicitly bounded
set of requested projections is complete. This gives the policy a defined
input and prevents a search-like task from becoming an unbounded bus or frame
latency source.

\subsection{\texttt{DECISION}: selection becomes an accountable boundary}

\texttt{DECISION} begins only after evaluation has assigned values to the
complete candidate set and \texttt{SELECT} has retained one candidate using a
deterministic ordering. It is therefore not another evaluation pass. Its
responsibility is to stabilize the selected move, preserve its diagnostic
trace, and make the result eligible for publication through the mailbox.

The outgoing transition from \texttt{DECISION}, rather than the state name
itself, writes the response data and sets \texttt{READY} last. This ordering
forms a small commit protocol: a reader never observes a ready result with
partially written move data. It also provides a precise audit point: every
external move can be associated with an accepted request sequence, a
rules-valid candidate, its projected trace, and the policy comparison that
selected it.

Deterministic tie-breaking is part of this boundary. When candidate values
are equal, a canonical ordering---for example origin square followed by
landing order---selects one move without hidden randomness. The result is
repeatable across host and target adapters, which is essential for fixtures,
replay, and comparison of future policy changes.

\subsection{Why the separation matters}

Without this split, a game program can blur three distinct questions: whether
a move is legal, what its consequences are, and whether it should be exposed
as the chosen action. The FSA assigns them to \texttt{GENERATE},
\texttt{ANTICIPATE}, and \texttt{SELECT}/\texttt{DECISION}, respectively.
That allocation does not guarantee strong play; it makes the source of a
choice testable. A future policy can become more sophisticated without
changing the validity of the rules, the bounded nature of projection, or the
mailbox commit discipline.

\subsection{Portable boundary}

The shared core is C++11 and avoids target-specific dependencies. Host
emulation and target firmware are adapters around a stable board, mailbox, and
agent interface. The intended target includes resource-constrained systems,
but hardware measurements are not yet reported.

\section{Validation methodology}

Validation proceeds in layers:
\begin{enumerate}
  \item Verify the 32-square board mapping, coordinate round trips, edge
  behavior, and invalid encodings with independent geometry fixtures.
  \item Validate mailbox request parsing, sequencing, and error responses.
  \item Fixture-test legal move generation and move application.
  \item Add deterministic policy tests over fixed board positions.
  \item Evaluate complete players against declared baselines using a
  documented, reproducible protocol.
\end{enumerate}

Only the first three layers are partially implemented in this draft. The
results section will remain intentionally incomplete until later layers and
their measurements exist.

\section{Current scope and limitations}

The current rules implementation covers uncrowned-men quiet moves and
single captures with compulsory-capture filtering. Promotion, king movement,
multi-jumps, draw rules, policy, deeper projection, target adapters, and
performance experiments remain future work. These limits are deliberate
boundaries, not claims of completed functionality.

\section{Related work and historical context}

This section is reserved for a reviewed comparison with Checkers programs,
finite-state game controllers, constrained game hardware, and period
computer-player implementations. The project has an initial lead list, but
citations will be added only after source verification. No novelty claim
should be inferred from this placeholder section.

\section{Evaluation plan}

The completed paper will report:
\begin{itemize}
  \item rules conformance against a documented position suite;
  \item determinism and protocol-error behavior;
  \item move quality against explicitly named baseline policies;
  \item memory, code size, and latency on each supported platform; and
  \item limitations, failure cases, and reproducibility instructions.
\end{itemize}

\section{Conclusion}

The project proposes a testable separation between game-control FSA, rules,
and policy for portable deterministic Checkers play. The contribution of this
draft is an accountable research direction and an implementation evidence
trail. Any stronger claim awaits the remaining rules, experiments, and
related-work review.

\end{document}
